\mainlanguage[de]

\setuplayout
  [backspace=2.0cm,
   height=middle,
   width=middle,
   leftmargin=1.8cm,
   footerdistance=0.0cm,
   rightmargindistance=0.2cm,
   rightmargin=1.8cm,
   bottom=0.5cm,
   leftmargindistance=0.2cm,
   header=2.1cm,
   headerdistance=0.5cm,
   footer=0.5cm,
   top=0cm,
   topspace=0.5cm,
   bottomspace=0.0cm]

\setupformulas[align=right]

\definehead[exercise][subsection]

\setuphead
  [exercise]
  [number=yes,
   style=,
   numberstyle=\ss\bf\tfb,
   textstyle=\ss,
   alternative=margin,
   numbercommand={\determineheadnumber[subsection]\currentheadnumber\gobbleoneargument}]

\setuphead
  [chapter]
  [number=no,
   style=,
   textstyle=\tfb\ss]

\setuphead
  [section]
  [number=no,
   style=,
   textstyle=\tf\ss]

\definemeasure[headerimageheight][\dimexpr(\headerheight-2ex)]

\definelayer[sabel]

\startsetups[sabel]

\setlayerframed[sabel][x=0mm,y=0mm][frame=off,offset=none]
{%
   \setupTABLE[start][style=\tfx]
   \setupTABLE[r][each][height=0.85\baselineskip]
   \setupTABLE[column][1][align=right, width=.4\textwidth, loffset=-2pt, toffset=-5pt]%
   \setupTABLE[column][2][align={center,lohi}, width=.2\textwidth]%
   \setupTABLE[column][3][align=left, width=.4\textwidth, roffset=-2pt, toffset=-5pt]%
   \bTABLE[frame=off]
   \bTR \bTD Sabel-GBS                     \eTD \bTD [nr=5] \dontleavehmode\externalfigure[assets/sabel.png][maxheight=\measure{headerimageheight}, frame=off]\dontleavehmode \eTD \bTD Semester WS 2025 \eTD \eTR
   \bTR \bTD Computergestützte Messtechnik \eTD                                                                                                                                     \bTD Klasse E3A       \eTD \eTR
   \bTR \bTD B.~Donner                     \eTD                                                                                                                                     \bTD                  \eTD \eTR
   \bTR \bTD \strut                        \eTD                                                                                                                                     \bTD                  \eTD \eTR
   \bTR \bTD 1.Schulaufgabe                \eTD                                                                                                                                     \bTD 17.12.2025       \eTD \eTR
   \eTABLE}

\stopsetups

\setupbackgrounds[header][text][bottomframe=on,setups=sabel,background=sabel]
\setuppagenumbering[location={footer,middle},style=\bfc]

\setupfootertexts
  [BE: 48 Arbeitszeit: 60min][Viel Erfolg!]
  [Viel Erfolg!][BE: 48 Arbeitszeit: 60min]


\definedescription[example][
  headstyle=\it, style=, align=left, location=top,
  width=broad, margin=1cm, indenting=never, text=\hskip-1cm]

\definemathcommand [kgV] [nolop] {\mfunction{kg\kern-.15em\hbox{V}}}

\setdigitmode 3

\setupitemgroup[itemize][1][n,repeat][width=1em]
\setupitemgroup[itemize][2][n,repeat][width=2em]
\setupitemgroup[itemize][3][n,repeat][width=3em]
\setupitemgroup[itemize][4][n,repeat][width=4em]

\startMPinclusions
  input "include/metapost/mymacros.mp"
\stopMPinclusions

\definefloat
   [combinationfloat]
   [combinationfloats]
   [figure]

\setupcaption
   [combinationfloat]
   [number=no]

\setupcaptions
   [location={top, left}]

\setupcombinations[location=top]

\starttext

\externalfigure[GBS-CMT/Datenblätter/Datenblatt_Hochdrucksensoren.pdf][width=\textwidth, frame=off]

\page[yes]

\externalfigure[GBS-CMT/Datenblätter/Datenblatt_Hochdrucksensoren.pdf][width=\textwidth, frame=off, page=2]

\page[yes]

Hinweise: Rechnen Sie auf drei signifikante Ziffern genau (drei
Nachkommastellen). Versehen Sie Zahlen mit den korrekten Einheiten in Ihrer
Lösung. Ersatzergebnisse geben nicht das korrekte Ergebnis wieder. Sie sind
nur zum Lösen der weiteren Teilaufgaben gedacht.


\startitemize
\item Mit folgendem Sensor soll der Druck in einer Bremsleitung eines
Fahrzeugs kontinuierlich überwacht werden. Die Versorgungsspannung des
Sensors beträgt in der ganzen Aufgabe \digits{5.000}\unit{Volt}.
\startitemize
\item Auf dem Datenblatt befinden sich die Angaben: Druckbereich, Berstdruck
und Überdruck max. Geben Sie die Grenzen der drei Druckbereiche an und
erläutern Sie kurz, was in den einzelnen Druckbereichen geschieht. (BE~6)
\item Berechnen Sie die Empfindlichkeit dieses Sensors laut Datenblatt. (BE~2)
\item Bestimmen Sie rechnerisch den Bereich der Ausgangsspannung am unteren
und am oberen Ende des Druckbereichs unter Berücksichtigung des Offset-- und
Empfindlichkeitsfehlers. Beachten Sie, dass sich der Offsetfehler laut
Datenblatt auf die Versorgungsspannung bezieht und der Empfindlichkeitsfehler
auf den gemessenen Druck. (BE~8)
\item Übertragen Sie die Kennline in ein korrekt und eindeutig beschriftetes
Diagramm. Markieren Sie entlang dieser Kennline das Fehlerband dass sich
insgesamt aus Offset-- und Empfindlichkeitsfehler ergibt. (1\unit{cm} pro
\unit{Mega Pascal} und \digits{2.5}\unit{cm} pro \unit{Volt})  (BE~6)
\item Bei einer Kalibrierung des Sensors haben Sie folgende Messwerte erhalten:

\startframedtext[frame=off,offset=overlay]
\starttabulate[|cw|ck1|ck1|ck1|ck1|ck1|ck1|][distance=none] 
    \FL
    \VL Druck in \unit{Mega Pascal}      \VL[3] \digits{0}    \VL  \digits{5}   \VL  \digits{10}  \VL  \digits{15}  \VL  \digits{20}  \VL  \digits{25}  \VL\AR
    \LL[3]
    \VL (Ist) Spannung in \unit{Volt}    \VL[3] \digits{0.59} \VL \digits{1.42} \VL \digits{2.14} \VL \digits{2.91} \VL \digits{3.86} \VL \digits{4.21} \VL\AR
    \LL[1]                                               
    \VL (Soll) Spannung in \unit{Volt}   \VL[3]               \VL               \VL               \VL               \VL               \VL               \VL\AR
    \LL[1]                                               
    \VL \Delta Spannung in \unit{Volt}   \VL[3]               \VL               \VL               \VL               \VL               \VL               \VL\AR
   \BL
\stoptabulate
\stopframedtext
\startitemize

\item[it:Aufg21] Bestimmen Sie eine optimale Kennline mithilfe der Methode der kleinsten
Quadrate. Geben Sie die Ausgleichsgerade für diese Messwerte auf drei gültige
Ziffern an.

[Ersatzergebnis zum Weiterrechnen, falls die Aufgabe nicht gelöst werden
konnte: \math{U_A(p) = \digits{0.170}\frac{\unit{Volt}}{\unit{Mega Pascal}}
\cdot p + \digits{0.651}\unit{Volt}}] (BE~4)

\item Die gefundene Kalibrierungskennline weicht von der Kennlinie aus dem
Datenblatt geringfügig ab. Bestimmen Sie den Offsetfehler bezogen auf die
Versorgungsspannung, den die Kalibrierungskennline gegenüber der
Datenblattkennlinie aufweist. (BE~3)

\item Bestimmen Sie den Empfindlichkeitsfehler bezogen auf den gemessen Druck,
den die Kalibrierungskennline gegenüber der Datenblattkennlinie aufweist.
(BE~3)

\item[it:Aufg24] Bestimmen Sie den verbleibenden Linearitätsfehler, bezüglich der
Kalibrierungskennline (optimale Kennline). Geben Sie den Linearitätsfehler in
Prozent--Fullscale an. (BE~6)

\page[yes]

\item Begründen Sie, ob dieser Sensor mit den obigen Messwerten noch innerhalb
der Datenblattangaben liegt. (BE~2)
\stopitemize

% \item In Teilaufgabe 1.4 haben Sie bereits den absoluten Fehler als Fehlerband
% um die Herstellerkennlinie eingezeichnet. Jetzt gilt es, den relativen Fehler
% sowohl am unteren als auch am oberen Ende des Druckbereichs zu bestimmen.
% Anschließend versuchen Sie bitte, den funktionalen Zusammenhang des relativen
% Fehlers in Bezug auf die Eingangsgröße Druck darzustellen. Das Ziel besteht
% darin, eine Funktion zu finden, die für jeden Eingangsdruck den entsprechenden
% relativen Fehler liefert. Mithilfe dieser Funktion könnte das relative
% Fehlerband grafisch dargestellt werden. (BE~4)

\item In der folgenden Aufgabe soll der ursprüngliche Drucksensor kalibriert
werden, allerdings haben Sie mit Einschränkungen zu kämpfen.

\startitemize

\item Sie dürfen aus Zeitgründen nur drei Messungen durchführen, um den
Sensor zu kalibrieren. Ihre Aufgabe ist es, zu erklären, wie Sie die drei
Messungen optimal aufteilen, um trotz der begrenzten Anzahl an Messungen die
bestmögliche Genauigkeit der Kalibrierung zu erreichen.

Hinweis: Der Sensor kann bei jedem Druckwert gleich wahrscheinlich von der
idealen (linearen) Kennlinie abweichen.
Begründen Sie Ihre Entscheidung für die Verteilung der Messpunkte. Gerne auch
mithilfe einer Skizze. (BE~2)

\item Leider hat Ihr Unternehmen Ihnen aus Kostengründen nur ein
unzureichendes Druckmessgerät zur Verfügung gestellt. Mit diesem Gerät können
Sie nur bei den Druckwerten 0, 5 und 10\unit{Mega Pascal} messen. Dabei
erhalten Sie die in der Tabelle angegebenen Spannungswerte.

Bleibt die Kalibrierungskennlinie im Vergleich zu \in{Teilaufgabe}[it:Aufg21]
unverändert? Begründen Sie Ihre Antwort entweder durch eine Rechnung oder
durch nachvollziehbare Argumente. (BE~2)

\item Begründen Sie, wie sich der absolute Linearitätsfehler im Vergleich zu
\in{Teilaufgabe}[it:Aufg24]
verändern wird. (BE~2)

\item Sie sollen einen Drucksensor mit einem Messbereich bis
25\unit{Mega Pascal} kalibrieren. Allerdings können Sie nur Vergleichsdaten
für die Druckwerte zwischen 0 und 10\unit{Mega Pascal} erheben.
Skizzieren Sie Messdaten, die mit einem solchen eingeschränkten Verfahren
niemals korrekt kalibriert werden können. Begründen Sie, warum diese
eingeschränkte Kalibrierung bei Ihren Messdaten scheitern muss. (BE~2)

\stopitemize
\stopitemize
\stopitemize

\stoptext
